
\documentclass{article}
\usepackage{colortbl}
\usepackage{makecell}
\usepackage{multirow}
\usepackage{supertabular}

\begin{document}

\newcounter{utterance}

\centering \large Interaction Transcript for game `cladder', experiment `full\_v1.5\_default', episode 9295 with selfrefine{-}llama{-}3.1{-}8b{-}instruct.
\vspace{24pt}

{ \footnotesize  \setcounter{utterance}{1}
\setlength{\tabcolsep}{0pt}
\begin{supertabular}{c@{$\;$}|p{.15\linewidth}@{}p{.15\linewidth}p{.15\linewidth}p{.15\linewidth}p{.15\linewidth}p{.15\linewidth}}
    \# & \multicolumn{2}{c}{Player} && \multicolumn{2}{c}{Game Master} \\
    \hline

    \theutterance \stepcounter{utterance}  
    & & & \multicolumn{4}{p{0.6\linewidth}}{
        \cellcolor[rgb]{0.9,0.9,0.9}{
            \makecell[{{p{\linewidth}}}]{
                \texttt{\tiny{[P1$\langle$GM]}}
                \texttt{You are an expert in causal inference. The following question is not a typical commonsense query, but rather a meticulously designed question created by a professor specializing in causal inference, intended to assess the students' mastery of the course content.} \\
\texttt{Imagine a self{-}contained, hypothetical world with only the following conditions, and without any unmentioned factors or causal relationships: Zory has a direct effect on zuph and glimx. Zuph has a direct effect on jyka. Jyka has a direct effect on glimx. Zory is unobserved. The overall probability of zuph is 96\%. The probability of not zuph and glimx is 3\%. The probability of zuph and glimx is 30\%. Is the chance of glimx smaller when observing zuph?} \\
\texttt{Start your answer with 'Yes' or 'No', followed by additional reasoning or evidence to support your explanation.} \\
            }
        }
    }
     \\ \\

    \theutterance \stepcounter{utterance}  
    & & \multicolumn{4}{p{0.6\linewidth}}{
        \cellcolor[rgb]{0.95,0.95,0.95}{
            \makecell[{{p{\linewidth}}}]{
                \texttt{\tiny{[P1$\leftrightarrow$P]}}
                \texttt{{[}DRAFT{]}} \\
\texttt{No.} \\
\\ 
\texttt{To address this question, we need to apply the principles of causal inference, specifically the concept of confounding variables and the backdoor criterion. } \\
\\ 
\texttt{Given the conditions:} \\
\\ 
\texttt{1. Zory has a direct effect on zuph and glimx.} \\
\texttt{2. Zuph has a direct effect on jyka.} \\
\texttt{3. Jyka has a direct effect on glimx.} \\
\texttt{4. Zory is unobserved.} \\
\\ 
\texttt{We can represent these relationships using a directed acyclic graph (DAG):} \\
\\ 
\texttt{Zory {-}> Zuph} \\
\texttt{Zory {-}> Glimx} \\
\texttt{Zuph {-}> Jyka} \\
\texttt{Jyka {-}> Glimx} \\
\\ 
\texttt{Since Zory is unobserved, we cannot directly measure its effect on Glimx. However, we can use the backdoor criterion to identify the causal effect of Zuph on Glimx.} \\
\\ 
\texttt{The backdoor criterion states that if we can block all backdoor paths from Zuph to Glimx, we can estimate the causal effect of Zuph on Glimx. In this case, the backdoor path is Zuph {-}> Jyka {-}> Glimx.} \\
\\ 
\texttt{To block this path, we need to condition on a set of variables that separates Zuph and Glimx. However, since Jyka is a direct effect of Zuph, conditioning on Jyka will not block the backdoor path.} \\
\\ 
\texttt{Now, let's examine the given probabilities:} \\
\\ 
\texttt{{-} P(Zuph) = 96\%} \\
\texttt{{-} P(¬Zuph ∧ Glimx) = 3\%} \\
\texttt{{-} P(Zuph ∧ Glimx) = 30\%} \\
\\ 
\texttt{We can use these probabilities to calculate the probability of Glimx given Zuph, P(Glimx|Zuph). Using Bayes' theorem, we can write:} \\
\\ 
\texttt{P(Glimx|Zuph) = P(Zuph ∧ Glimx) / P(Zuph)} \\
\texttt{= 30\% / 96\%} \\
\texttt{= 0.3125} \\
\\ 
\texttt{Now, let's calculate the probability of Glimx given ¬Zuph, P(Glimx|¬Zuph). We can use the law of total probability:} \\
\\ 
\texttt{P(Glimx|¬Zuph) = P(¬Zuph ∧ Glimx) / P(¬Zuph)} \\
\texttt{= 3\% / (1 {-} 96\%)} \\
\texttt{= 3\% / 4\%} \\
\texttt{= 0.075} \\
\\ 
\texttt{Comparing the two probabilities, we see that P(Glimx|Zuph) > P(Glimx|¬Zuph). This suggests that the presence of Zuph increases the probability of Glimx.} \\
\\ 
\texttt{However, we need to consider the effect of Zory, which is unobserved. Since Zory has a direct effect on both Zuph and Glimx, it may be confounding the relationship between Zuph and Glimx} \\
            }
        }
    }
    & \\ \\

    \theutterance \stepcounter{utterance}  
    & & \multicolumn{4}{p{0.6\linewidth}}{
        \cellcolor[rgb]{0.95,0.95,0.95}{
            \makecell[{{p{\linewidth}}}]{
                \texttt{\tiny{[P1$\leftrightarrow$P]}}
                \texttt{{[}SELF{-}FEEDBACK{]}} \\
\texttt{FORMAT: ok} \\
\texttt{RULES: ok} \\
\texttt{STRATEGY: problem} \\
\texttt{REASON: The answer incorrectly concludes that the presence of Zuph increases the probability of Glimx, when in fact the question asks if the chance of Glimx is smaller when observing Zuph.} \\
            }
        }
    }
    & \\ \\

    \theutterance \stepcounter{utterance}  
    & \multicolumn{4}{p{0.6\linewidth}}{
        \cellcolor[rgb]{0.561,0.737,0.561}{
            \makecell[{{p{\linewidth}}}]{
                \texttt{\tiny{[P1$\rangle$GM]}}
                \texttt{No.} \\
\\ 
\texttt{The presence of Zuph does not necessarily imply a smaller chance of Glimx. In fact, the probability of Glimx given Zuph, P(Glimx|Zuph), is 0.3125, which is greater than the probability of Glimx given ¬Zuph, P(Glimx|¬Zuph), which is 0.075. However, this does not directly answer the question. To determine if the chance of Glimx is smaller when observing Zuph, we need to consider the effect of Zory, which is unobserved. Since Zory has a direct effect on both Zuph and Glimx, it may be confounding the relationship between Zuph and Glimx.} \\
            }
        }
    }
    & & \\ \\

    \theutterance \stepcounter{utterance}  
    & & & \multicolumn{2}{p{0.3\linewidth}}{
        \cellcolor[rgb]{0.9,0.9,0.9}{
            \makecell[{{p{\linewidth}}}]{
                \texttt{\tiny{[GM$|$GM]}}
                \texttt{No} \\
            }
        }
    }
    & & \\ \\

    \theutterance \stepcounter{utterance}  
    & & & \multicolumn{2}{p{0.3\linewidth}}{
        \cellcolor[rgb]{0.9,0.9,0.9}{
            \makecell[{{p{\linewidth}}}]{
                \texttt{\tiny{[GM$|$GM]}}
                \texttt{yes} \\
            }
        }
    }
    & & \\ \\

    \theutterance \stepcounter{utterance}  
    & & & \multicolumn{2}{p{0.3\linewidth}}{
        \cellcolor[rgb]{0.9,0.9,0.9}{
            \makecell[{{p{\linewidth}}}]{
                \texttt{\tiny{[GM$|$GM]}}
                \texttt{game\_result = LOSE} \\
            }
        }
    }
    & & \\ \\

\end{supertabular}
}

\end{document}
